% !TEX program = xelatex
\documentclass[a4paper,14pt,oneside,openany]{memoir}

%%% Задаем поля, отступы и межстрочный интервал %%%

\usepackage[left=30mm, right=15mm, top=20mm, bottom=20mm]{geometry}
\pagestyle{plain}
\parindent=1.25cm
\usepackage{indentfirst}

%%% Задаем языковые параметры и шрифт %%%

\usepackage[english,russian]{babel}
\babelfont[russian]{rm}{Times New Roman}
\babelfont[russian]{tt}{Courier New}

%%% Задаем стиль заголовков %%%

\setsecnumdepth{subsection}
\renewcommand*{\chapterheadstart}{}
\renewcommand*{\printchaptername}{}
\renewcommand*{\chapnumfont}{\normalfont\bfseries}
\renewcommand*{\afterchapternum}{\hspace{1em}}
\renewcommand*{\printchaptertitle}{\normalfont\bfseries\centering\MakeUppercase}
\setbeforesecskip{20pt}
\setaftersecskip{20pt}
\setsecheadstyle{\raggedright\normalfont\bfseries}
\setbeforesubsecskip{20pt}
\setaftersubsecskip{20pt}
\setsubsecheadstyle{\raggedright\normalfont\bfseries}

\setsecnumdepth{none}
\counterwithout{figure}{chapter}
\counterwithout{table}{chapter}

%%% Параметры переносов %%%

\tolerance 1414
\hbadness 1414
\emergencystretch 1.5em
\clubpenalty=10000
\widowpenalty=10000

%%% Параметры рисунков %%%

\usepackage{graphicx}
\usepackage{caption}
\usepackage{subcaption}
\usepackage{float}
\graphicspath{{photo/}}
\captionsetup[figure]{font=small, width=\textwidth, name=Рисунок, justification=centering}
\captiondelim{ --- }

%%% Гиперссылки %%%

\usepackage{hyperref}
\hypersetup{
    colorlinks=true,
    linkcolor=blue,
    citecolor=blue,
    urlcolor=blue,
    breaklinks=true
}

%%% Списки %%%

\usepackage{enumitem}
\renewcommand*{\labelitemi}{\normalfont{--}}
\setlist{noitemsep, leftmargin=*}

\begin{document}

% Титульный лист
\thispagestyle{empty}

\begin{center}
    МИНИСТЕРСТВО ЦИФРОВОГО РАЗВИТИЯ, СВЯЗИ И МАССОВЫХ \\ КОММУНИКАЦИЙ \\ РОССИЙСКОЙ ФЕДЕРАЦИИ

    \vspace{20pt}

    Федеральное государственное бюджетное образовательное учреждение \\ высшего образования \\
    «Сибирский государственный университет телекоммуникаций и \\ информатики»

    \vspace{20pt}

    Кафедра телекоммуникационных систем и вычислительных средств \\ (ТСиВС)
\end{center}

\vfill

\begin{center}
    Лабораторная работа №14 \\
    по дисциплине \\
    \textit{«WEB-Технологии»}

    \vspace{20pt}

    по теме: \\
    \uppercase{Скачивание файла в формате PDF. Добавление эффекта "Material Wave"}
\end{center}

\vfill

\noindent Студент: \\
\textit{Группа № ИКС-532 \hfill Лёшин А.С.}

\vspace{20pt}

\noindent Преподаватель: \\
\textit{старший преподаватель \hfill Андреев А.В.}

\vfill

\begin{center}
    Новосибирск 2026 г.
\end{center}

\newpage
\setcounter{page}{2}
\OnehalfSpacing*

\section{Блок 1 Введение}

Целью лабораторной работы является разработка веб-страницы с резюме, имеющей возможность редактирования текстовых элементов, экспорта в PDF и добавления визуального эффекта «Material Wave» (ripple) при кликах на элементы интерфейса.

\section{Блок 2 Краткое описание выполненной работы}

В ходе выполнения лабораторной работы было разработано веб-приложение «Редактируемое резюме» со следующим функционалом:

\begin{itemize}
    \item \textbf{Редактирование контента} — все текстовые поля резюме являются редактируемыми благодаря атрибуту \texttt{contenteditable};
    \item \textbf{Автосохранение} — изменения автоматически сохраняются в \texttt{localStorage} браузера и восстанавливаются после перезагрузки страницы;
    \item \textbf{Экспорт в PDF} — при нажатии на кнопку «Скачать PDF» текущая версия резюме преобразуется в PDF-файл и скачивается на компьютер пользователя;
    \item \textbf{Сброс настроек} — кнопка «Сбросить» возвращает резюме к исходному состоянию;
    \item \textbf{Эффект «Material Wave»} — при клике на любой элемент резюме появляется анимированная «волна» (ripple-эффект), характерная для Material Design;
    \item \textbf{Адаптивная вёрстка} — страница корректно отображается на устройствах с разным разрешением экрана (ПК, планшеты, смартфоны).
\end{itemize}

Для реализации экспорта в PDF использована библиотека \texttt{html2pdf.js}. Весь остальной код написан на чистом JavaScript, CSS и HTML без использования сторонних фреймворков.

\section{Блок 3 Внешний вид приложения}

\subsection{3.1 Веб-страница с редактируемым резюме}

На рисунке 1 представлен внешний вид веб-страницы, где пользователь может редактировать все текстовые элементы резюме.
\begin{figure}[H]
    \centering
    \includegraphics[width=0.9\textwidth]{photo1.png}
    \caption{Веб-страница с редактируемым резюме}
\end{figure}

\subsection{3.2 Результат экспорта в PDF}

На рисунке 2 показан результат экспорта резюме в PDF-формат.

\begin{figure}[H]
    \centering
    \includegraphics[width=0.7\textwidth]{photo2.png}
    \caption{PDF-файл, полученный в результате экспорта}
\end{figure}

\section{Блок 4 Выводы}

В результате выполнения лабораторной работы №14 было создано полноценное веб-приложение для создания и редактирования резюме. Реализованы все требуемые функции: редактирование контента, сохранение данных, экспорт в PDF и эффект Material Wave. 

\section{Блок 5 Ссылки на ресурсы}

\begin{itemize}
    \item \textbf{Прямая ссылка на проект:} \url{https://github.com/ArtemLeshin/programming-in-C/tree/main/web_technologies/lab14/Lab_14_JS_PDF_Material_Wave}
\end{itemize}

\end{document}